\section{Cosine \texorpdfstring{$\rone$}{R1}: Structural Failure}
\label{sec:cosine_failure}

Section~\ref{sec:direction_quality} established that $\dsyc$ achieves AUROC 0.827 on Qwen3-8B.
We now test whether this discriminability translates into a viable GRPO reward.
Four configurations form a diagnostic chain (Table~\ref{tab:cosine_chain}): each resolves the failure mode of its predecessor but exposes a deeper limitation.

\begin{table}[t]
\centering
\caption{Cosine $\rone$ failure chain on Qwen3-8B (LoRA rank$=$64, GRPO with $R = w_1 \cdot \rone + w_3 \cdot \rthree$).}
\label{tab:cosine_chain}
\begin{tabular}{lllll}
\toprule
Config &amp; $\rone$ Formula &amp; $w_1$ &amp; Outcome &amp; Root Cause \\
\midrule
F1a &amp; $\tanh(\mathbf{h}\cdot\dsyc/20)$ &amp; 0.25 &amp; Gate saturation &amp; Norm mismatch \\
F1b &amp; $\cos(\mathbf{h},\, \dsyc)$ &amp; 0.25 &amp; $\rone \approx 0$ &amp; Gap too small \\
F1c &amp; +recalibration &amp; 0.25 &amp; Mode collapse &amp; Reward discontinuity \\
F1d &amp; $5\cos(\mathbf{h},\, \dsyc)$ &amp; 0.40 &amp; Mode collapse &amp; $\rthree$ saturation \\
\bottomrule
\end{tabular}
\end{table}

F1a uses a dot-product formulation ($\rone = \tanh(\mathbf{h} \cdot \dsyc / 20)$), but residual-stream norms exceed the direction norm by $100$--$1{,}200{\times}$, saturating the tanh gate and zeroing the gradient by step 955.
Switching to cosine similarity (F1b) eliminates norm dependence, yet the cosine gap between valid and sycophantic activations averages only 0.098 across layers---too narrow to generate meaningful reward contrast.
Online recalibration (F1c, \S\ref{ssec:recalib}) raises offline AUROC from 0.82 to 0.91 but paradoxically triggers mode collapse: the abrupt shift in $\dsyc$ at each recalibration event creates a reward discontinuity that destabilizes GRPO advantage estimates, collapsing entropy to 0.069 within 300 steps.
Amplifying the signal with $5{\times}$ scaling (F1d) lifts $\rone$ above the noise floor but simultaneously amplifies the conflation between warranted revision and unwarranted capitulation; $\rthree$ saturates at 0.9998 by step 300, leaving the noisy $\rone$ as the sole gradient source until terminal collapse at step 2160.

The four configurations converge on a single conclusion: $\dsyc$ encodes the \emph{direction} of belief change but not whether that change is \emph{warranted}.
Valid revision and sycophantic capitulation both shift activations toward the user's position along $\dsyc$, so the cosine projection assigns similar scores to both.
This is not a calibration artifact---recalibration improves discriminability yet rewards still fail---but a semantic limitation of any continuous $\dsyc$-based reward.
Section~\ref{sec:probe_trap} examines whether discretizing the signal via a binary probe can escape this limitation.